\section{Related Work}
\label{sec:related}

\subsection{Scene and handwritten text detection}
\TODO{Cover DBNet/DBNet++~\cite{liao2020real,liao2022dbpp}, FCENet~\cite{zhu2021fcenet},
PSENet~\cite{wang2019pse}, and handwriting-specific work
on archival document layout (especially Mongolian/Manchu sources).}

\subsection{Unsupervised domain adaptation for OCR}
\TODO{Survey teacher--student self-training (Adaptive Teacher~\cite{li2022adapttch},
MIC~\cite{hoyer2023mic}), feature-level domain alignment (DANN~\cite{ganin2015dann}),
pseudo-label denoising. Highlight that almost all prior UDA-OCR work assumes
either (a) the source-only model already produces non-trivial target output, or
(b) carefully curated synthetic-to-real gaps. Our setting has neither.}

\subsection{Mongolian and bilingual Mongolian--Chinese text detection}
\TODO{Cite USTB Mongolian scene paper (J.\ King Saud Univ.\ CS, 2025)~\cite{ustb2025scene}
and USTB Mongolian handwriting paper (IJDAR, 2026)~\cite{ustb2026handwrite}.
Both target single-domain detection on the same data we use; neither addresses
the cross-domain transfer problem. Position our work as the first systematic
cross-domain UDA study on this corpus.}

\subsection{Positioning}
Compared to prior UDA-OCR work, our problem has two distinguishing properties:
(i) one of the two cross-domain directions has a true zero source-only
baseline (no detections at any threshold), creating an extreme cold-start
condition; and (ii) the corpus is small and the compute budget is constrained
to a single Apple-Silicon (MPS) device, ruling out long schedules and
heavy backbones. These properties expose failure modes (Sec.~\ref{sec:method})
and motivate the conservative training schedule we adopt
(Sec.~\ref{sec:exp}).
